
\documentclass{article}
\usepackage{colortbl}
\usepackage{makecell}
\usepackage{multirow}
\usepackage{supertabular}

\begin{document}

\newcounter{utterance}

\centering \large Interaction Transcript for game `cladder', experiment `full\_v1.5\_default', episode 8837 with reflexion{-}Qwen3.5{-}9B.
\vspace{24pt}

{ \footnotesize  \setcounter{utterance}{1}
\setlength{\tabcolsep}{0pt}
\begin{supertabular}{c@{$\;$}|p{.15\linewidth}@{}p{.15\linewidth}p{.15\linewidth}p{.15\linewidth}p{.15\linewidth}p{.15\linewidth}}
    \# & \multicolumn{2}{c}{Player} && \multicolumn{2}{c}{Game Master} \\
    \hline

    \theutterance \stepcounter{utterance}  
    & & & \multicolumn{4}{p{0.6\linewidth}}{
        \cellcolor[rgb]{0.9,0.9,0.9}{
            \makecell[{{p{\linewidth}}}]{
                \texttt{\tiny{[P1$\langle$GM]}}
                \texttt{You have attempted this same task before and did not succeed. These are the plans you wrote after those attempts.} \\
\texttt{Trial \#1: New plan: I failed to complete the task because I stopped my reasoning prematurely after calculating the Natural Indirect Effect (NIE) and did not calculate the Natural Direct Effect (NDE) or the Total Effect, which are necessary to determine if the effect through the mediator is negative. Specifically, I ignored the provided probabilities for \$P(M=1 | Y=1, Z=0)\$ and \$P(M=1 | Y=1, Z=1)\$, which are required to compute the NDE (\$P(M=1|Y=1, Z=0) {-} P(M=1|Y=1, Z=1)\$). Next time, I will explicitly calculate the NDE using the given values (\$57\textbackslash{}\% {-} 10\textbackslash{}\% = +47\textbackslash{}\%\$), compare it with the NIE (\$+7\textbackslash{}\%\$), and conclude that since both direct and indirect effects are positive, the answer is "No" (it does not negatively affect), ensuring I utilize all provided conditional probabilities before finalizing the response.} \\
\\ 
\\ 
\texttt{You are an expert in causal inference. The following question is not a typical commonsense query, but rather a meticulously designed question created by a professor specializing in causal inference, intended to assess the students' mastery of the course content.} \\
\texttt{Imagine a self{-}contained, hypothetical world with only the following conditions, and without any unmentioned factors or causal relationships: Yupt has a direct effect on zupj and muvq. Xyfo has a direct effect on zupj and muvq. Zupj has a direct effect on muvq. Xyfo is unobserved. For those who are not yupt and are not zupj, the probability of muvq is 44\%. For those who are not yupt and are zupj, the probability of muvq is 51\%. For those who are yupt and are not zupj, the probability of muvq is 57\%. For those who are yupt and are zupj, the probability of muvq is 10\%. For those who are not yupt and are not xyfo, the probability of zupj is 42\%. For those who are not yupt and are xyfo, the probability of zupj is 87\%. For those who are yupt and are not xyfo, the probability of zupj is 26\%. For those who are yupt and are xyfo, the probability of zupj is 44\%. The overall probability of xyfo is 22\%. Does yupt negatively affect muvq through zupj?} \\
\texttt{Start your answer with 'Yes' or 'No', followed by additional reasoning or evidence to support your explanation.} \\
            }
        }
    }
     \\ \\

    \theutterance \stepcounter{utterance}  
    & \multicolumn{4}{p{0.6\linewidth}}{
        \cellcolor[rgb]{0.561,0.737,0.561}{
            \makecell[{{p{\linewidth}}}]{
                \texttt{\tiny{[P1$\rangle$GM]}}
                \texttt{No} \\
\\ 
\texttt{**Reasoning:**} \\
\\ 
\texttt{To determine if \$Y\_\{upt\}\$ negatively affects \$M\_\{uvq\}\$ through \$Z\_\{upj\}\$, we must calculate the **Natural Indirect Effect (NIE)**. The NIE represents the effect of changing the mediator \$Z\_\{upj\}\$ from its value under the control condition (\$Y=0\$) to its value under the treatment condition (\$Y=1\$), while holding the treatment \$Y\$ fixed at the control level (\$Y=0\$).} \\
\\ 
\texttt{The formula for the NIE is:} \\
\texttt{\$\$NIE = P(M=1 | Y=1, Z=Z\_\{Y=1\}) {-} P(M=1 | Y=0, Z=Z\_\{Y=1\})\$\$} \\
\texttt{In simpler terms, it is the difference in the probability of the outcome \$M\$ when the mediator \$Z\$ is set to the value it would take if \$Y=1\$, compared to when \$Z\$ is set to the value it would take if \$Y=0\$, keeping \$Y\$ fixed at 0.} \\
\\ 
\texttt{**Step 1: Identify the probabilities of the mediator \$Z\$ under different conditions of \$Y\$.**} \\
\texttt{We are given the following probabilities for \$Z\_\{upj\}\$:} \\
\texttt{*   \$P(Z=1 | Y=0) = 42\textbackslash{}\%\$ (For those not \$Y\_\{upt\}\$)} \\
\texttt{*   \$P(Z=1 | Y=1) = 26\textbackslash{}\%\$ (For those who are \$Y\_\{upt\}\$ and not \$X\_\{yfo\}\$) and \$44\textbackslash{}\%\$ (For those who are \$Y\_\{upt\}\$ and \$X\_\{yfo\}\$).} \\
\\ 
\texttt{*Correction/Refinement on Mediator Distribution:*} \\
\texttt{The problem states \$X\_\{yfo\}\$ is unobserved. However, the probabilities for \$Z\$ given \$Y\$ are provided conditionally on \$X\$ as well. To find the marginal \$P(Z=1 | Y=1)\$, we must average over the distribution of \$X\$.} \\
\texttt{Given:} \\
\texttt{*   \$P(X=1) = 22\textbackslash{}\%\$} \\
\texttt{*   \$P(X=0) = 78\textbackslash{}\%\$} \\
\texttt{*   \$P(Z=1 | Y=1, X=0) = 26\textbackslash{}\%\$} \\
\texttt{*   \$P(Z=1 | Y=1, X=1) = 44\textbackslash{}\%\$} \\
\\ 
\texttt{Calculate \$P(Z=1 | Y=1)\$:} \\
\texttt{\$\$P(Z=1 | Y=1) = P(Z=1 | Y=1, X=0)P(X=0) + P(Z=1 | Y=1, X=1)P(X=1)\$\$} \\
\texttt{\$\$P(Z=1 | Y=1) = (0.26 \textbackslash{}times 0.78) + (} \\
            }
        }
    }
    & & \\ \\

    \theutterance \stepcounter{utterance}  
    & & & \multicolumn{2}{p{0.3\linewidth}}{
        \cellcolor[rgb]{0.9,0.9,0.9}{
            \makecell[{{p{\linewidth}}}]{
                \texttt{\tiny{[GM$|$GM]}}
                \texttt{No} \\
            }
        }
    }
    & & \\ \\

    \theutterance \stepcounter{utterance}  
    & & & \multicolumn{2}{p{0.3\linewidth}}{
        \cellcolor[rgb]{0.9,0.9,0.9}{
            \makecell[{{p{\linewidth}}}]{
                \texttt{\tiny{[GM$|$GM]}}
                \texttt{yes} \\
            }
        }
    }
    & & \\ \\

    \theutterance \stepcounter{utterance}  
    & & & \multicolumn{2}{p{0.3\linewidth}}{
        \cellcolor[rgb]{0.9,0.9,0.9}{
            \makecell[{{p{\linewidth}}}]{
                \texttt{\tiny{[GM$|$GM]}}
                \texttt{game\_result = LOSE} \\
            }
        }
    }
    & & \\ \\

    \theutterance \stepcounter{utterance}  
    & & \multicolumn{4}{p{0.6\linewidth}}{
        \cellcolor[rgb]{0.95,0.95,0.95}{
            \makecell[{{p{\linewidth}}}]{
                \texttt{\tiny{[P1$\leftrightarrow$P]}}
                \texttt{Post{-}episode Reflexion (not sent to GM):} \\
\texttt{New plan: I failed to complete the task because I stopped my reasoning prematurely after calculating the Natural Indirect Effect (NIE) and did not calculate the Natural Direct Effect (NDE) or the Total Effect, which are necessary to determine if the effect through the mediator is negative. Specifically, I ignored the provided probabilities for \$P(M=1 | Y=1, Z=0)\$ and \$P(M=1 | Y=1, Z=1)\$, which are required to compute the NDE (\$P(M=1|Y=1, Z=0) {-} P(M=1|Y=1, Z=1)\$). Next time, I will explicitly calculate the NDE using the given values (\$57\textbackslash{}\% {-} 10\textbackslash{}\% = +47\textbackslash{}\%\$), compare it with the NIE (\$+7\textbackslash{}\%\$), and conclude that since both direct and indirect effects are positive, the answer is "No" (it does not negatively affect), ensuring I utilize all provided conditional probabilities before finalizing the response.} \\
            }
        }
    }
    & \\ \\

\end{supertabular}
}

\end{document}
